% GENERATED FILE. Edit canonical lessons, then run npm run generate:books.
% canonical-source-hash: fnv1a64:8803030ce90602b0
% canonical-lessons: AR-C364-hallala, AR-C364-unsur, AR-C364-miyar, AR-C364-haddada, AR-C364-mayyaza

\chapter{Analysing things}
\label{ch:ar-analysing-things}
\hlchaptermodality{364}

\begin{chapteropening}
\textbf{By the end of this chapter:} \emph{I can say that someone analysed, specified or told things apart, and talk about elements and criteria.}

Complete the last lesson of chapter 364: I can say that someone analysed, specified or told things apart, and talk about elements and criteria.
\end{chapteropening}

\section[\texorpdfstring{\ar{حلل} (ḥallala)}{حلل (ḥallala)}]{\ar{حلل} (ḥallala) — to analyse}

\label{lesson:AR-C364-hallala}



\begin{quote}
Before the new one: say the Arabic for a phenomenon, then the Arabic for to interpret.
\end{quote}

\subsection*{You'll want to know: \ar{حلل}}
\textbf{\ar{حلل}} — \emph{ḥallala} — \textquotedblleft{}to analyse\textquotedblright{}.

\begin{quote}
\ar{حلل} \ar{الطالب} \ar{القصيدة} \ar{بعناية}.
\end{quote}

\begin{quote}
\emph{ḥallala aṭ-ṭālib al-qaṣīda bi-ʿināya.}
\end{quote}

\begin{quote}
The student analysed the poem carefully.
\end{quote}

Present: \ar{يحلل}, \emph{yuḥallil}.

\begin{cousinweb}
\textbf{\ar{حلل}} (\emph{ḥallala}) is built on the native Arabic root \emph{ḥ-l-l}, \textquotedblleft{}to untie, to loosen\textquotedblright{} (Lane). It is Form II; in modern use it means to undo a thing into its parts: to analyse.
\end{cousinweb}

\begin{grammarlens}[title={What you've built}]
One more word for giving reasons and opinions.
\end{grammarlens}

\subsection*{Your turn}
\begin{itemize}
\raggedright
  \item \emph{Say it:} \emph{ḥallala}
  \item \emph{Say it:} \emph{ḥallala}, once more
  \item \emph{Say it:} \emph{ḥallala aṭ-ṭālib al-qaṣīda bi-ʿināya.} — The student analysed the poem carefully.
  \item \emph{Recall:} say the Arabic for a phenomenon, then the Arabic for to interpret, then say \emph{ḥallala} again
\end{itemize}

\begin{tcolorbox}[breakable,colback=teal!4,colframe=teal!35!black,title={Before you move on}]
What does \ar{حلل} mean? (\textquotedblleft{}To analyse\textquotedblright{}.) Say it once more.
\end{tcolorbox}
\section[\texorpdfstring{\ar{عنصر} (ʿunṣur)}{عنصر (ʿunṣur)}]{\ar{عنصر} (ʿunṣur) — an element}

\label{lesson:AR-C364-unsur}



\begin{quote}
Before the new one: say the Arabic for to interpret, then the Arabic for to analyse.
\end{quote}

\subsection*{You'll want to know: \ar{عنصر}}
\textbf{\ar{عنصر}} — \emph{ʿunṣur} — \textquotedblleft{}an element\textquotedblright{}.

\begin{quote}
\ar{الماء} \ar{عنصر} \ar{مهم} \ar{في} \ar{الطبخ}.
\end{quote}

\begin{quote}
\emph{al-māʾ ʿunṣur muhimm fī aṭ-ṭabkh.}
\end{quote}

\begin{quote}
Water is an important element in cooking.
\end{quote}

Plural \ar{عناصر}, \emph{ʿanāṣir}.

\begin{cousinweb}
\textbf{\ar{عنصر}} (\emph{ʿunṣur}) is an old word for the stock or basic matter of a thing, which Lane uses when he explains other words. Its origin is not settled.
\end{cousinweb}

\begin{grammarlens}[title={What you've built}]
One more word for giving reasons and opinions.
\end{grammarlens}

\subsection*{Your turn}
\begin{itemize}
\raggedright
  \item \emph{Say it:} \emph{ʿunṣur}
  \item \emph{Say it:} \emph{ʿunṣur}, once more
  \item \emph{Say it:} \emph{al-māʾ ʿunṣur muhimm fī aṭ-ṭabkh.} — Water is an important element in cooking.
  \item \emph{Recall:} say the Arabic for to interpret, then the Arabic for to analyse, then say \emph{ʿunṣur} again
\end{itemize}

\begin{tcolorbox}[breakable,colback=teal!4,colframe=teal!35!black,title={Before you move on}]
What does \ar{عنصر} mean? (\textquotedblleft{}An element\textquotedblright{}.) Say it once more.
\end{tcolorbox}
\section[\texorpdfstring{\ar{معيار} (miʿyār)}{معيار (miʿyār)}]{\ar{معيار} (miʿyār) — a criterion}

\label{lesson:AR-C364-miyar}



\begin{quote}
Before the new one: say the Arabic for to analyse, then the Arabic for an element.
\end{quote}

\subsection*{You'll want to know: \ar{معيار}}
\textbf{\ar{معيار}} — \emph{miʿyār} — \textquotedblleft{}a criterion\textquotedblright{}.

\begin{quote}
\ar{الجودة} \ar{هي} \ar{المعيار} \ar{الأول} \ar{عندنا}.
\end{quote}

\begin{quote}
\emph{al-jawda hiya al-miʿyār al-awwal ʿindanā.}
\end{quote}

\begin{quote}
Quality is our first criterion.
\end{quote}

Plural: \ar{معايير}, \emph{maʿāyīr}.

\begin{cousinweb}
\textbf{\ar{معيار}} (\emph{miʿyār}) is built on the native Arabic root \emph{ʿ-y-r} (Lane). It is a noun of instrument, usually traced to the checking of weights and measures: a standard that others are measured against.
\end{cousinweb}

\begin{grammarlens}[title={What you've built}]
One more word for giving reasons and opinions.
\end{grammarlens}

\subsection*{Your turn}
\begin{itemize}
\raggedright
  \item \emph{Say it:} \emph{miʿyār}
  \item \emph{Say it:} \emph{miʿyār}, once more
  \item \emph{Say it:} \emph{al-jawda hiya al-miʿyār al-awwal ʿindanā.} — Quality is our first criterion.
  \item \emph{Recall:} say the Arabic for to analyse, then the Arabic for an element, then say \emph{miʿyār} again
\end{itemize}

\begin{tcolorbox}[breakable,colback=teal!4,colframe=teal!35!black,title={Before you move on}]
What does \ar{معيار} mean? (\textquotedblleft{}A criterion\textquotedblright{}.) Say it once more.
\end{tcolorbox}
\section[\texorpdfstring{\ar{حدد} (ḥaddada)}{حدد (ḥaddada)}]{\ar{حدد} (ḥaddada) — to specify}

\label{lesson:AR-C364-haddada}



\begin{quote}
Before the new one: say the Arabic for an element, then the Arabic for a criterion.
\end{quote}

\subsection*{You'll want to know: \ar{حدد}}
\textbf{\ar{حدد}} — \emph{ḥaddada} — \textquotedblleft{}to specify\textquotedblright{}.

\begin{quote}
\ar{حدد} \ar{المدير} \ar{وقت} \ar{الاجتماع}.
\end{quote}

\begin{quote}
\emph{ḥaddada al-mudīr waqt al-ijtimāʿ.}
\end{quote}

\begin{quote}
The manager set the time of the meeting.
\end{quote}

Present: \ar{يحدد}, \emph{yuḥaddid}.

\ar{حدود}, borders, comes from the same root.

\begin{cousinweb}
\textbf{\ar{حدد}} (\emph{ḥaddada}) is built on the native Arabic root \emph{ḥ-d-d}, \textquotedblleft{}to sharpen; to set a limit\textquotedblright{} (Lane). It is Form II, to mark the limits of: to define, to fix, to specify.
\end{cousinweb}

\begin{grammarlens}[title={What you've built}]
One more word for giving reasons and opinions.
\end{grammarlens}

\subsection*{Your turn}
\begin{itemize}
\raggedright
  \item \emph{Say it:} \emph{ḥaddada}
  \item \emph{Say it:} \emph{ḥaddada}, once more
  \item \emph{Say it:} \emph{ḥaddada al-mudīr waqt al-ijtimāʿ.} — The manager set the time of the meeting.
  \item \emph{Recall:} say the Arabic for an element, then the Arabic for a criterion, then say \emph{ḥaddada} again
\end{itemize}

\begin{tcolorbox}[breakable,colback=teal!4,colframe=teal!35!black,title={Before you move on}]
What does \ar{حدد} mean? (\textquotedblleft{}To specify\textquotedblright{}.) Say it once more.
\end{tcolorbox}
\section[\texorpdfstring{\ar{ميز} (mayyaza)}{ميز (mayyaza)}]{\ar{ميز} (mayyaza) — to distinguish}

\label{lesson:AR-C364-mayyaza}



\begin{quote}
Before the new one: say the Arabic for a criterion, then the Arabic for to specify.
\end{quote}

\subsection*{You'll want to know: \ar{ميز}}
\textbf{\ar{ميز}} — \emph{mayyaza} — \textquotedblleft{}to distinguish\textquotedblright{}.

\begin{quote}
\ar{لا} \ar{يميز} \ar{الطفل} \ar{بين} \ar{الألوان} \ar{بعد}.
\end{quote}

\begin{quote}
\emph{lā yumayyizu aṭ-ṭifl bayna al-alwān baʿd.}
\end{quote}

\begin{quote}
The child cannot tell the colours apart yet.
\end{quote}

Present: \ar{يميز}, \emph{yumayyiz}.

Often with \ar{بين}: \ar{ميز} \ar{بين} … \ar{و} …, told … from ….

\begin{cousinweb}
\textbf{\ar{ميز}} (\emph{mayyaza}) is built on the native Arabic root \emph{m-y-z}, \textquotedblleft{}to separate, to set apart\textquotedblright{} (Lane). It is Form II, to set apart: to tell one thing from another.
\end{cousinweb}

\begin{grammarlens}[title={What you've built}]
One more word for giving reasons and opinions.
\end{grammarlens}

\subsection*{Your turn}
\begin{itemize}
\raggedright
  \item \emph{Say it:} \emph{mayyaza}
  \item \emph{Say it:} \emph{mayyaza}, once more
  \item \emph{Say it:} \emph{lā yumayyizu aṭ-ṭifl bayna al-alwān baʿd.} — The child cannot tell the colours apart yet.
  \item \emph{Recall:} say the Arabic for a criterion, then the Arabic for to specify, then say \emph{mayyaza} again
\end{itemize}

\begin{tcolorbox}[breakable,colback=teal!4,colframe=teal!35!black,title={Before you move on}]
What does \ar{ميز} mean? (\textquotedblleft{}To distinguish\textquotedblright{}.) Say it once more.
\end{tcolorbox}
